\subsection{Gemeinsame Stetige Verteilung}
\shortdefinition $\cX_i$ haben gemeinsame stetige Dichte wenn Abbildung $f : \R^n \rightarrow \R_+$ existert, so dass $\forall a_i \in \R$ gilt:
\[
    \P[\cX_1 \! \leq \! a_1, \ldots, \cX_n \! \leq \! a_n]\! = \! \int_{-\8}^{a_1}\! \ldots \! \int_{-\8}^{a_n} \! f(x_1, \ldots) \dx x_n \ldots \mbox{d} x_1
\]

\shorttheorem $f$ gem. Dichte
\[
    \int_{-\8}^{\8}\! \ldots \! \int_{-\8}^{\8} \! f(x_1, \ldots) \dx x_n \ldots \mbox{d} x_1 = 1
\]
Jeder Funk. $f$ die obiges erfüllt ein W-Raum $(\Omega, \cF, \P)$ und $n$-Z.V. $\cX_i$ zugeo. werden, s.d. $f$ die gem. Dichte von $\cX_i$ ist.

\inlineintuition $f(x_1, \ldots) \dx x_1, \ldots$ beschreibt die W., dass ein zu-fälliger Punkt $(\cX_1, \ldots)$ in $[x_1, x_1 + \dx x_1] \times \ldots$ liegt

\shortexample Gleichverteilungen:
\begin{itemize}
    \item \bi{Einheitsquadrat} $f(x, y) = \begin{cases}
                  1 & (x, y) \in [0, 1]^2 \\
                  0 & \text{sonst}
              \end{cases}$
    \item \bi{Einheitskreisscheibe} $D = \{ (x, y) \in \R^2 \divider x^2 + y^2 \leq 1 \}$
          $f(x, y) = \begin{cases}
                  \frac{1}{\pi} & (x, y) \in D \\
                  0             & \text{sonst} \\
              \end{cases}$
\end{itemize}

\shorttheorem[Erwartungswert Bild]
\[
    \E[\varphi(\cX_1, \ldots)] \! = \! \int_{-\8}^{\8} \ldots \int_{-\8}^{\8} \varphi(x_1, \ldots) f(x_1, \ldots) \dx x_n \ldots
\]
{\scriptsize $\varphi$ ist bspw. für $\E[\cX\cY]$, $\varphi(\cX, \cY) = \cX \cdot \cY$, Integralgrenzen nach $f(x_1, \ldots)$}

\shortdefinition[Randverteilung] Falls $\cX, \cY$ gemeinsame Verteilung $F$ haben, so ist die Verteilungsfunktion der \bi{Randverteilung} von $\cX$,
$F_\cX : \R \rightarrow [0, 1]$ gegeben durch:
\[
    x \mapsto F_\cX(x) \! = \! \P[\cX \leq x] \! = \! \P[\cX \leq x, \cY < \8] \! = \! \limit{y}{\8} F(x, y)
\]
Analog für $\cY$ ist sie $F_y : \R \rightarrow [0, 1]$:
\[
    y \mapsto F_\cY(y) \! = \! \P[\cY \leq y] \! = \! \P[\cY < \8, \cY \leq y] \! = \! \limit{x}{\8} F(x, y)
\]
Die Dichten der Randverteilungen sind:
\[
    f_\cX(x) = \int_{-\8}^{\8} f(x, y) \dx y
    \qquad
    f_\cY(y) = \int_{-\8}^{\8} f(x, y) \dx x
\]

\shortexample Randdichten mit Gleichverteilung\\
\bi{Einheitsquadrat} {\scriptsize ($f(x, y) = 1_{(x, y) \in [0, 1]^2} = 1_{x \in [0, 1]} 1_{y \in [0, 1]}$)}
\begin{align*}
    f_\cX(x) & = \int_{0}^{1} 1_{x \in [0, 1]} 1_{y \in [0, 1]} \dx y = 1_{x \in [0, 1]} \\
    f_\cY(y) & = \int_{0}^{1} 1_{x \in [0, 1]} 1_{y \in [0, 1]} \dx x = 1_{y \in [0, 1]}
\end{align*}

\bi{Einheitskreisscheibe} {\scriptsize ($f(x, y) = \frac{1}{\pi} \cdot 1_{x^2 + y^2 \leq 1}$)}
\begin{align*}
    f_\cX(x) \! = \! \int_{-\8}^{\8} \frac{1}{\pi} 1_{y^2 \leq 1 - x^2} \dx y \! = \! \int_{-\sqrt{1 - x^2}}^{\sqrt{1 - x^2}} \frac{1}{\pi} \dx y \!
    = \! \frac{2}{\pi} \sqrt{1 - x^2} \\
    f_\cY(y) \! = \! \int_{-\8}^{\8} \frac{1}{\pi} 1_{x^2 \leq 1 - y^2} \dx x \! = \! \int_{-\sqrt{1 - y^2}}^{\sqrt{1 - y^2}} \frac{1}{\pi} \dx x \! = \! \frac{2}{\pi} \sqrt{1 - y^2}
\end{align*}

\shorttheorem[Unabhängigkeit] $\cX_i$ unabhängig g.d.w.
\[
    f_{\cX_1, \ldots, \cX_n}(x_1, \ldots, x_n) = f_{\cX_1}(x_1) \cdot \ldots \cdot f_{\cX_n}(x_n)
\]

\shortexample Unabhängigkeit bei Gleichverteilungen\\
\bi{Einheitsquadrat} Koordinaten sind unabhängig

\bi{Einheitskreisscheibe} Koordinaten nicht unabhängig
